\section{Introduction and Motivation}
Loop structures are ubiquitous in programming and form the foundation of iterative reasoning. In traditional programs such as C, determining suitable invariants for \texttt{for} or \texttt{while} loops remains difficult, often relying on heuristics or manual insight. Similar iterative behaviors also appear in AI agents that repeatedly observe their surroundings, make decisions based on the current state, and iterate over actions or conditions until a specific goal is achieved~\cite{evertz2003loop, wu2022survey, mosqueira2023human}. Ensuring the correctness of these program-level loops is therefore crucial for the reliable functioning of such systems.  
% Loop structures are ubiquitous in programming, including in embodied AI agents. For example, an embodied agent may repeatedly observe its surroundings and make decisions based on the current state of the environment. Such behaviors often involve loops that iterate over actions or conditions until a specific goal is achieved

A well-established approach to reasoning about loop correctness is through \emph{loop invariants}—properties that remain true before and after each loop iteration~\cite{furia2014loop, d2008survey}. Over the years, several methods have been proposed to automatically generate loop invariants.

% The first category includes \textbf{traditional methods}, such as constraint-based techniques
Traditional approaches—including constraint solving~\cite{Coln2003LinearIG,8603013}, counterexample-guided refinement~\cite{clarke2000counterexample,bogomolov2017counterexample}, and Craig interpolation~\cite{jhala2006practical, mcmillan2010lazy} —encode invariants as logical constraints and iteratively refine them through symbolic reasoning. Although effective on small programs, these techniques require domain expertise to design templates or logical rules and often fail to scale to real-world software, where each new data structure demands specialized inference patterns.

Machine learning–based approaches leverage large code corpora to learn syntactic and semantic patterns and automatically generate invariants~\cite{garg2014ice,garg2016learning,nguyen2017counterexample,si2020code2inv}. These data-driven techniques generally offer better scalability than symbolic ones, as they reduce dependence on manually designed templates and rules by learning invariant patterns from data. 
%These methods are generally more scalable and efficient than traditional ones. 
However, they often suffer from issues such as overfitting, limited interpretability, and a strong dependence on the quality and diversity of the training data, leading to degraded performance on unseen programs or data structures.
%. Moreover, their performance often degrades significantly when applied to entirely new or previously unseen datasets.

%The most recent trend involves \textbf{large language model (LLM)–based methods}, which use LLMs to infer loop invariants from source code and prompts
More recently, large language model (LLM)–based approaches infer loop invariants directly from program code and natural-language prompts~\cite{wu2024llm,cao2025clause2inv,LIU2026103387}. These approaches typically follow a \emph{guess-and-check} paradigm: the LLM generates candidate invariants that are subsequently verified using SMT solvers or formal reasoning tools, with counterexamples fed back for refinement. 
%the LLM first generates candidate invariants, which are then verified using formal reasoning tools such as SMT solvers. When an invariant fails verification, feedback from the verifier is used to iteratively refine the generated invariants. 
Such methods are generally more flexible and adaptable, leveraging the extensive knowledge and contextual understanding embedded in LLMs.

However, despite their impressive capabilities, current LLMs still struggle to produce correct and precise invariants for programs with complex control flows, such as conditional branches. This limitation arises partly from the transformer architecture, which, while powerful in capturing global context, lacks explicit mechanisms for fine-grained structural reasoning and long-range code dependencies.

To address this challenge, we propose \thistool, a novel approach that combines the precision of program analysis that captures branch correlations with the generative and contextual strengths of LLMs. Our approach is designed to improve the precision and robustness of invariant generation for complex loops. This direction holds promise for advancing automated reasoning in program verification and synthesis tasks.
